% TOP_PROBLEM: {"id": 1500006, "title": "Algebraic Stories — The $k$-rank of monomials", "category_id": 4, "category": {"id": 4, "name": "algebra", "display_name": "Algebra", "description": "Group theory, ring theory, field theory, and algebraic structures.", "slug": "algebra", "order_index": 4, "created_at": "2026-07-31T15:26:25.671Z"}, "status": "partially_solved"}
% FIRST_POSTED: null
% ENTRY_KIND: statement_only
% Imported from: batch06.tex; UnsolvedMath ID 1500006
% Compile twice: lualatex 1500006.tex
\documentclass[11pt]{article}
\usepackage{fontspec}
\setmainfont{Latin Modern Roman}
\usepackage{amsmath,amssymb,mathtools,unicode-math}
\setmathfont{Latin Modern Math}
\usepackage{xurl,hyperref}
\hypersetup{colorlinks=true,urlcolor=blue,pdftitle={UnsolvedMath Problem 1500006}}
\newcommand{\sref}[2]{\hyperlink{source:#1:#2}{[S#2]}}
\newcommand{\cataloguescope}[1]{\paragraph{Mathematical question.}#1}
\begin{document}
% UnsolvedMath ID 1500006; source catalogue code AMR-014-0006
\setcounter{section}{1500005}
\section{Algebraic Stories — The $k$-rank of monomials}\label{1500006}
{\small\sffamily UnsolvedMath ID 1500006 · Algebra}
\subsection{Definitions and mathematical statement}

\paragraph{Shared notation.}$\mathbb N=\{1,2,\ldots\}$, and $\mathbb Z,\mathbb Q,\mathbb R,\mathbb C$ denote the integers, rationals, reals and complex numbers. Any other conventions are specified locally.


Let $S=\mathbb C[x,y]$, and let $S_e$ denote its homogeneous degree-$e$ forms. For $f\in S_{kd}$, let $\operatorname{rk}_k(f)$ be the smallest number of terms in a decomposition $f=\sum_jg_j^k$ with $g_j\in S_d$. Define $R_{k,d}=\max_{f\in S_{kd}}\operatorname{rk}_k(f)$, the maximal $k$-rank of binary forms.Let $k\geq3$, $d\geq1$ and $a,b\geq0$ satisfy $a+b=kd$. Define $s,t\in\{0,\ldots,k-1\}$ by $a\equiv s\pmod k$ and $b\equiv t\pmod k$.
\paragraph{Statement recorded in the supplied source.}
The bound $\operatorname{rk}_k(x^ay^b)\leq\max(s,t)+1$ is given in the source. Is it always an equality?
\[\operatorname{rk}_k(x^ay^b)=\max(s,t)+1.\] \sref{1500006}{2}
\subsection{Short English statement}
Is the known upper bound for binary monomials, determined by their exponents modulo $k$, always exact?
\subsection{Sources}
\begin{description}
\item[\hypertarget{source:1500006:1}{S1}] ULAM.AI and the UnsolvedMath Contributors, UnsolvedMath: A Curated Collection of Open Mathematics Problems, record 1500006 (AMR-014-0006). CC BY 4.0. \url{https://huggingface.co/datasets/ulamai/UnsolvedMath}.
\item[\hypertarget{source:1500006:2}{S2}] R. Fröberg, S. Lundqvist, A. Oneto and B. Shapiro, \emph{Algebraic Stories from One and from the Other Pockets}, arXiv:1801.01692 (2018), §1.3, Problem D. \url{https://arxiv.org/abs/1801.01692}.
\end{description}
\label{end:1500006}
\end{document}
